\chapter{Two Ways a Concept Can Be Fixed}
\label{ch:two-ways}

\begin{plainlanguage}
\textbf{In plain terms:} Is intelligence something real we're discovering, or something we're collectively inventing through argument? This chapter stakes out a middle position: part of what "intelligence" means is forced on us by reality and won't budge; part is just convention and could have gone differently. The tools to tell which part is which come later.
\end{plainlanguage}

Before building any machinery, it is worth being precise about what success would even look like --- what it would mean for the elimination process to work, and what kind of thing the survivor would be once it had worked. There are two very different answers available, and the difference between them is not a technicality. It determines whether this entire project is a contribution to epistemology (we are getting closer to a truth that was already there) or to something closer to sociology of concepts (we are watching a category get built, with no further fact about what it ``really'' is once the building is done).

\textbf{Discovery.} On this view, there is a real, mind-independent fact about what intelligence is --- a $\theta^*$ somewhere in $\Theta$ --- and the elimination process, whatever form it takes, is a method for approximating $\theta^*$ more and more closely. Mistakes are possible (a definition can be wrongly eliminated, or a bad one wrongly retained), and the measure of the process's success is how close the eventually-converged $\Theta_T$ comes to containing or approximating $\theta^*$.

\textbf{Constitution.} On this view, there is no $\theta^*$ waiting to be found. The word ``intelligence'' does not refer to a pre-existing natural kind; it refers to whatever survives the elimination process, full stop. $\Theta_T$ does not approximate anything external to itself --- it \emph{is} the thing the word picks out, because the elimination process is not measurement, it is concept-formation. There is no further question to ask, on this view, about whether $\Theta_T$ got it right, because ``getting it right'' presupposes exactly the fixed external target that constitution denies exists.

These positions are not merely different emphases. They make different predictions about what should happen if the elimination process were rerun from a different starting point, by a different culture, with a different sequence of perspectives encountered. Discovery predicts convergence to (approximately) the same $\Theta_T$ regardless of path, because there is a real target pulling every honest inquiry toward it. Constitution predicts no such guarantee --- a differently-ordered, differently-populated elimination process could land somewhere else entirely, with no sense in which one outcome is more correct than the other.

\section*{The working position}

This book does not adopt either pole. It adopts what will be called \emph{weak realism}, and stating it precisely now --- even before the machinery exists to justify it --- sets the target for everything that follows.

\begin{definition}[Hard/Soft Decomposition]
Let
\[
\Theta_t = H_t \cup S_t
\]
where $H_t$ (the \emph{hard core}) is the portion of the current candidate set fixed by content that has survived contact with resistant, non-negotiable reality --- a system's measurable incapacity, a hazard that reliably produces the outcome it threatens, a counterexample that is simply, checkably true --- and $S_t$ (the \emph{soft periphery}) is the portion fixed only by local, revisable, social agreement: emphasis, caricature, rhetorical amplification, cultural priority.
\end{definition}

\begin{conjecture}[Differential Stability]
As evidence accumulates, $dH_t/dt \to 0$ --- the hard core stabilizes and stops moving, because what fixes it (contact with reality that doesn't change its mind) doesn't go away. $dS_t/dt$, by contrast, may remain nonzero indefinitely --- the soft periphery can keep shifting forever, because nothing forces it to stop.
\end{conjecture}

This is stated as a conjecture rather than a theorem because nothing in this chapter yet defines what ``evidence accumulation'' formally does to $H$ and $S$ --- that requires the elimination machinery of Chapter~\ref{ch:fidelity}, in particular the fidelity functional that will let the book distinguish, for any given candidate constraint, whether it belongs to the kind of evidence that stabilizes $H$ or the kind that merely reshuffles $S$. Until that machinery exists, the hard/soft split is an organizing intuition, not a derived result. It is the intuition the rest of the book is built to earn.

If the conjecture holds, weak realism follows naturally: intelligence is neither purely discovered (there is no single fixed $\theta^*$, only a hard core that itself may be multidimensional rather than a single point) nor purely constituted (the hard core is not up for grabs --- no amount of cultural rearrangement will make a system that provably cannot do $X$ count as having done $X$). The interesting work, on this view, is not deciding the fork once and for all. It is learning, case by case, which parts of any given definition belong to $H$ and which belong to $S$ --- which is exactly the question the elimination machinery of Part~II is designed to answer, and exactly the question that, by the book's own logic, the book's own argument about intelligence must itself survive being asked of.
